\section{실험 환경}
\label{sec:실험환경}

제안하는 \method{} 및 통합 프레임워크의 라우팅 유효성을 정밀하게 검증하기 위해 구성한 시뮬레이션 물리 파라미터 환경과 대조용 비교 알고리즘 기준선(Baselines)들을 기술한다.

\subsection{시뮬레이션 환경 및 물리 파라미터}
성능 평가의 공정성과 재현성을 위해 자체 구축된 로컬 Lite-GLOBE 시뮬레이터 환경을 사용하였다. 
시뮬레이션 가상 공간은 $1000\text{m} \times 1000\text{m} \times 1000\text{m}$ 크기의 3차원 공역으로 설정되었으며, 고정 육상 시설물인 소스 노드는 $[0\text{m}, 0\text{m}, 0\text{m}]$ 좌표에, 목적지 노드는 $[1000\text{m}, 1000\text{m}, 0\text{m}]$에 독립적으로 위치한다. 이 두 지상 기지국 사이의 중계를 담당하는 무인 항공기(UAV)들은 총 8대, 16대, 24대 등으로 구성을 스위핑하였으며, 고도 주행 제한 한계는 $[200\text{m}, 1000\text{m}]$ 범주로 세팅되었다. 각 무인기는 3차원 공간 상에서 최소 $-30\text{m}$에서 최대 $30\text{m}$의 축별 무작위 3차원 변위를 보이는 고속의 Random Waypoint 이동성 모델에 의거하여 기동하며, 무선 통신 유효 거리는 $\delta = 550\text{m}$로 고정되었다. 하나의 에피소드는 총 $T_{\mathrm{ep}} = 100$ 타임스텝의 고정 주기를 가지며, 신뢰성 평가를 위한 마지막 20타임스텝 동안의 패킷 생존 및 전달 성공 유지를 최종 메트릭(Hit Ratio)으로 합산한다.

주요 시뮬레이션 매개변수의 상세 구성은 표~\ref{tab:sim_params_kci}에 요약되어 있다.

\begin{table}[htbp]
\centering
\caption{Lite-GLOBE 시뮬레이션 매개변수 설정 요약}
\label{tab:sim_params_kci}
\begin{tabular}{lc}
\toprule
매개변수명 & 설정값 \\
\midrule
시뮬레이션 공역 크기 & $1000\text{m} \times 1000\text{m} \times 1000\text{m}$ \\
소스 및 목적지 좌표 & $[0,0,0]$ / $[1000,1000,0]$ \\
배치 UAV 노드 수 ($N$) & 8, 16, 24대 \\
UAV 통신 반경 ($\delta$) & 550m \\
UAV 고도 제한 범위 & $200\text{m} \sim 1000\text{m}$ \\
UAV 기동 제한 변위 범위 & $[-30\text{m}, 30\text{m}]$ (축별 단일 홉) \\
에피소드 최대 길이 ($T_{\mathrm{ep}}$) & 100 timesteps \\
성능 평가 관찰 윈도우 ($\epsilon$) & 20 timesteps \\
\bottomrule
\end{tabular}
\end{table}

\subsection{평가 시나리오 구성}
네트워크 구조적 장애 요인과 다이내믹 토폴로지 변화에 대응하는 능력을 고차원적으로 평가하기 위해 총 14가지 세부 평가 시나리오를 정의하였다.
\begin{itemize}
    \item \textbf{Nominal Scenarios}: 정규 밀도 및 단순 노이즈 링크 상태.
    \item \textbf{Routing-Hole Scenarios}: 지상 전파 차단이나 산악 장애물 배치 등으로 인해 목적지 직선 방향 노드 연결이 차단되어, GPSR의 로컬 미니멈 오류 유도가 불가피한 구조적 장애 토폴로지.
    \item \textbf{Severe Link-Loss Scenarios}: 기후 변화나 심각한 무선 다중 경로 페이딩으로 인해 패킷 손실 확률 $p_{\mathrm{loss}}$가 최대 30\%에 도달하는 극한 채널 상태.
    \item \textbf{Mobility/Scale stress}: UAV 대수가 급변하거나 노드 이동 속도가 정상 속도 대비 2배 이상 가속되는 과적 부하 시나리오.
\end{itemize}

\subsection{비교 대상 알고리즘 (Baselines)}
제안 기법인 \method{} 및 전신인 \phaseTwelve{} 모델의 상대적 가치를 파악하기 위해 다음 다섯 가지 기준선과 비교를 전개하였다.
\begin{enumerate}
    \item \textbf{GPSR} \cite{gpsr}: 순수 물리 거리 기반 Greedy 지리적 라우팅 대표 모델.
    \item \textbf{Predictive Geographic} (Heuristic): 1-hop 칼만 링크 잔여 예측치를 가중하여 포워딩하는 수학적 휴리스틱 모델.
    \item \textbf{Evo-QGeo} \cite{evoqgeo}: Q학습에 미래 상태 천이 예측 구조를 결합한 대표적인 강화학습 기반 지리 라우팅.
    \item \textbf{IQMR Q($\lambda$)} \cite{iqmr}: 다중홉 피드백 보상을 반영하여 우회 경로를 수렴시키는 Q-학습 기법.
    \item \textbf{DRAMA} \cite{drama}: 이웃 간의 실시간 학습 메신저(Emergent communication)를 320 Byte 크기로 연속 브로드캐스트하는 분산 CTDE MARL 프로토콜.
    \item \textbf{Student (No Switch)}: 위험 판단 게이트를 탑재하지 않고, 상시 MLP 신경망 추론과 예측만으로 차홉 결정을 진행하는 상시 구동형 대조군.
\end{enumerate}
